\chapter{Conclusão} \label{cha:conclusao}

Este projeto propõe uma solução analítica cientificamente fundamentada e aplicável para instituições financeiras que buscam ingressar no mercado de crédito rural de forma eficiente e segura, mesmo sem possuir bases de dados históricas. Ao utilizar os microdados regulatórios públicos do SICOR como fundamento analítico primário, o projeto transforma um contexto de assimetria de informação em um ambiente de decisão estruturada, demonstrando que a inteligência de dados públicos pode ser um diferencial competitivo legítimo e de alto valor estratégico.

Os impactos técnicos esperados do projeto incluem:
\begin{itemize}
    \item[(i)] Construção de um pipeline analítico reprodutível e documentado para o tratamento dos dados do SICOR, passível de manutenção e atualização à medida que novas safras agrícolas sejam incorporadas à base;
    \item[(ii)] Panorama estatístico detalhado da inadimplência por atividade, região, safra e finalidade, provendo à instituição uma visão robusta do perfil de risco dos diferentes subsegmentos do agronegócio;
    \item[(iii)] Modelo preditivo de risco de crédito rural agnóstico, validado por métricas padronizadas e interpretáveis, capaz de suportar a originação inicial de crédito rural na ausência de portfólio históricos próprio;
    \item[(iv)] Dashboard operacional de BI que traduz os insights em visualizações acessíveis e acionáveis para as equipes de crédito e de risco da instituição.
\end{itemize}

Do ponto de vista da metodologia científica, o projeto demonstra a viabilidade de construir modelos preditivos de risco de crédito a partir de dados públicos e regulatórios — uma abordagem com ampla aplicabilidade para outros segmentos financeiros onde a escassez de dados históricos constitui uma barreira de entrada relevante. A validação empírica das hipóteses formuladas contribuirá para o conhecimento a respeito do tema, especialmente em relação à capacidade de generalização dos modelos preditivos construídos a partir de dados públicos.

A convergência entre o ecossistema de BI e a modelagem preditiva de risco, ancorada em dador reais e acessíveis publicamente, posiciona este projeto como um caso de referência replicável para estratégias \textit{data-driven} no setor financeiro, contribuindo simultaneamente para o sucesso comercial da instituição, maturação das práticas de analytics aplicado a este setor e avanço do conhecimento científico em gestão de risco de crédito rural.
